\SourceLeaf{36}

\Needspace{8\baselineskip}
\begin{center}\small
BEMERKUNGEN ZU DEM AUFSATZE »BEWEIS DER IRRATIONALITÄT DER TANGENTEN
RATIONALER BÖGEN IN EINER NEUEN GESTALT«.
\end{center}
\par\noindent\hspace*{\parindent}Der erste einwandfreie Beweis, dass $\operatorname{tg}\dfrac{m}{n}$ für ganze Zahlen
$m, n$ irrational ist, wurde von \textsc{Legendre} in der vierten Anmerkung seiner »Éléments
de géométrie« (1794) entwickelt. Die Ausführungen von \textsc{Gauss} stimmen in den
Grundgedanken mit dem Beweis von \textsc{Legendre} überein. Neu ist die Art, wie
\textsc{Gauss} direkt aus den Potenzreihen der $P_\theta$ die unbegrenzte Abnahme dieser
Größen bei wachsendem $\theta$ nachweist. Über die hier einschlägigen
Untersuchungen \textsc{Lamberts} und \textsc{Legendres}, insbesondere über die vom
Letzteren an\,der\,zitierten Stelle vorgelegte Darstellung hat \textsc{Gauss} eine kritische
Besprechung aufgezeichnet, die hier jedoch nicht abgedruckt ist.

Eine andere Entwicklung, welche gleichfalls nicht abgedruckt ist, leitet aus den
Rekursionsformeln der $P_\theta$ angenäherte Darstellungen von $\operatorname{tg} m$
durch Quotienten ganzer rationaler Funktionen von $m$ ab. \textsc{Gauss} verwendet diese
Darstellungen zur Berechnung von Näherungswerten für $\tfrac{1}{2}\pi$.

\begin{flushright}
\textsc{Fricke}.
\end{flushright}
